\enquestion{4.51}{
    Assume that $Z$ has a standard normal distribution.
    Determine the value for $z$ that solves each of the following:
    
    \begin{tabular}{*{2}{m{.45\textwidth}}}
        (a) $P(Z<z)=0.9$ & (b) $P(Z<z)=0.5$\\
        (c) $P(Z > z)=0.1$ & (d) $P(Z > z)=0.9$ \\
        (e) $P(-1.28 < Z < z)=0.8$ & \\
    \end{tabular}
}
\enanswer{
    We can determine these boundary values $z$ using the graphical calculator using either the function \enquote{invnorm} (using Tail: Left if not specified) or the \enquote{numeric solver}:
    \begin{enumerate}[label=(\alph*)]
        \item $P(Z < z) = 0.9$ holds for $z = \invnorm(\opp=0.9, \mu=0, \sigma=1) \approx 1.2816$.
        \item $P(Z < z) = 0.5$ holds for $z = \invnorm(\opp=0.5, \mu=0, \sigma=1) = 0$ (notice that the standard normal distribution is symmetric in $x=0$).
        \item $P(Z > z) = 0.1$ holds for $z = \invnorm(\opp=0.1, \mu=0, \sigma=1, \text{Tail}=\text{Right}) \approx 1.2816$.
        \item $P(Z > z) = 0.9$ holds for $z = \invnorm(\opp=0.9, \mu=0, \sigma=1, \text{Tail}=\text{Right}) \approx -1.2816$.
        \item We can rewrite this equation as follows:
        \begin{align*}  
            P(Z < z) - P(Z < -1.28) &= 0.8 \Rightarrow P(Z < z) = 0.8 + P(Z < -1.28)
        \end{align*}
        This gives us
        \begin{align*}
            P(Z < z) &= 0.8 + \normalcdf(\lb=-\GRinfinity, \ub=-1.28, \mu=0, \sigma=1) \approx 0.9003 \\
                    &\Rightarrow z = \invnorm(\opp=0.9003, \mu=0,\sigma=1) \approx 1.2831.            
        \end{align*}
    \end{enumerate}
}

\enquestion{4.55}{
    Assume that $X$ is normally distributed with a mean of $5$ and a standard deviation of $4$.
    Determine the following:
    
    \begin{tabular}{*{3}{m{.3\textwidth}}}
        (a) $P(X<11)$ & (b) $P(X>2)$ & (c) $P(3<X<7)$\\
        (d) $P(-2<X<9)$ & (e) $P(2<X<8)$ & \\
    \end{tabular}
}
\enanswer{
    The following probabilities can be calculated using the function \enquote{normalcdf} on the graphical calculator:
    \begin{enumerate}[label=(\alph*)]
        \item $P(X < 11) = \normalcdf(\lb=-\GRinfinity, \ub=11, \mu=5, \sigma=4) \approx 0.9332$.
        \item $P(X > 2) = \normalcdf(\lb=2, \ub=\GRinfinity, \mu=5, \sigma=4) \approx 0.7734$.
        \item $P(3 < X < 7) = \normalcdf(\lb=3, \ub=7, \mu=5, \sigma=4) \approx 0.3829$.
        \item $P(-2 < X < 9) = \normalcdf(\lb=-2, \ub=9, \mu=5, \sigma=4) \approx 0.8013$.
        \item $P(2 < X < 8) = \normalcdf(\lb=2, \ub=8, \mu=5, \sigma=4) \approx 0.5467$.
    \end{enumerate}
}

\enquestion{4.56}{
    Assume that $X$ is normally distributed with a mean of $5$ and a standard deviation of $4$.
    Determine the value for $x$ that solves each of the following:

    \begin{tabular}{*{2}{m{.45\textwidth}}}
        (a) $P(X>x)=0.5$ & (b) $P(X>x)=0.95$ \\
        (c) $P(x<X<7)=0.2$ & (d) $P(3<X<x)=0.95$ \\
        (e) $P(-x<X-5<x)=0.99$ & \\
    \end{tabular}
}
\enanswer{
    We can determine these boundary values $x$ using the graphical calculator using either the function \enquote{invnorm} (using Tail: Left if not specified) or the \enquote{numeric solver}:
    \begin{enumerate}[label=(\alph*)]
        \item $P(X > x) = 0.5$ holds for $x$ equal to the mean, so $x = 5$.
        \item $P(X > x) = 0.95$ holds for $x = \invnorm(\opp=0.95, \mu=5, \sigma=4, \text{Tail}=\text{Right}) \approx -1.5794$.
        \item We can rewrite this equation as follows:
        \begin{align*}  
            P(X < 7) - P(X < x) = 0.2   &\Rightarrow P(X < x) = P(X < 7) - 0.2.
        \end{align*}
        This gives us
        \begin{align*}
            P(X < x)    &= \normalcdf(\lb=-\GRinfinity, \ub=7, \mu=5, \sigma=4) - 0.2 \approx 0.4915. \\
                        &\Rightarrow z = \invnorm(\opp=0.4915, \mu=5,\sigma=4) \approx 4.9144.   
        \end{align*}         
        \item We can rewrite this equation as follows:
        \begin{align*}  
            P(X < x) - P(X < 3) = 0.95  &\Rightarrow P(X < x) = 0.95 + P(X < 3).
        \end{align*}
        This gives us
        \begin{align*}
            P(X < x) = 0.95 + \normalcdf(\lb=-\GRinfinity, \ub=3, \mu=5, \sigma=4) \approx 1.2585.
        \end{align*}
        Notice that this equation is ill-defined, as a probability cannot be larger than one.
        \item $P(Z > z) = 0.9$ holds for $z = \invnorm(\opp=0.9, \mu=0, \sigma=1, \text{Tail}=\text{Right}) \approx -1.2816$.
        \item Notice that $X - 5$ is normally distributed with mean $0$ and variance $4$.
        In that case, $-x$ should be at the $0.5$-th percentile of the distribution and $x$ at the $99.5$-th percentile:
        \[
            x = \invnorm(\opp=0.995, \mu=0, \sigma=4) \approx 10.3033.
        \]
    \end{enumerate}
}

\enquestion{4.63}{
    The fill volume of an automated filling machine used for filling cans of carbonated beverage is normally distributed with a mean of $x$ cc and a standard deviation of $5$ cc.
    \begin{enumerate}[label=(\alph*)]
        \item At what value should the mean be set so that \SI{99.9}{\percent} of all cans exceed $365$ cc?
        \item At what value should the mean be set so that \SI{99.9}{\percent} of all cans exceed $365$ cc if the standard deviation can be reduced to $2$ cc?
    \end{enumerate}
}
\enanswer{
    The fill volume $X$ is normally distributed with mean $x$ and standard deviation $5$.
    \begin{enumerate}[label=(\alph*)]
        \item We want to solve for $x$ the equation $P(X > 365) = 0.999$.
        As the mean is unknown, we need to use the numeric solver for this:
        \begin{align*}
            E_1 &= \normalcdf(\lb=365, \ub=\GRinfinity, \mu=X, \sigma=5) \\
            E_2 &= 0.995
        \end{align*}
        The numeric solver gives us a mean $x \approx 380.4513$.
        \item We want to solve for $x$ the equation $P(X > 365) = 0.999$.
        As the mean is unknown, we need to use the numeric solver for this:
        \begin{align*}
            E_1 &= \normalcdf(\lb=365, \ub=\GRinfinity, \mu=X, \sigma=2) \\
            E_2 &= 0.995
        \end{align*}
        The numeric solver gives us a mean $x \approx 371.1805$.
    \end{enumerate}
}

\enquestion{4.69}{
    The life of a semiconductor laser at a constant power is normally distributed with a mean of $7000$ hours and a standard deviation of $600$ hours.
    \begin{enumerate}[label=(\alph*)]
        \item What is the probability that a laser fails before $6000$ hours?
        \item What is the life in hours that \SI{95}{\percent} of the lasers exceed?
        \item If three lasers are used in a product and they are assumed to fail independently, what is the probability that all three are still operating after $7000$ hours?
    \end{enumerate}
}
\enanswer{
    \begin{enumerate}[label=(\alph*)]
        \item The probability that a laser fails before $6000$ hours is equal to
        \begin{align*}
            P(X < 6000) &= \normalcdf(\lb=-\GRinfinity, \ub=6000, \mu=7000, \sigma=600) \\
            &\approx 0.0478.
        \end{align*}
        \item The life in hours that \SI{95}{\percent} of the lasers exceed is the value $x$ for which $P(X > x) = 0.95$.
        We can find $x$ by using the \enquote{invnorm} function:
        \[
            x = \invnorm(\opp=0.95, \mu=7000, \sigma=600, \text{Tail}=\text{Right}) \approx 6013
        \]
        \item The probability that a single laser still operates after $7000$ hours is equal to
        \[
            P(X > 7000) = \normalcdf(\lb=7000, \ub=\GRinfinity, \mu=7000, \sigma=600) = 0.5.
        \]
        Then, if we assume all three lasers fail independently, the probability that all three still operates after $7000$ hours is $0.5 \cdot 0.5 \cdot 0.5 = 0.125$.
    \end{enumerate}
}

\enquestion{4.91}{
    Suppose $X$ has an exponential distribution with mean equal to $10$.
    Determine the following:

    \begin{tabular}{*{3}{m{.3\textwidth}}}
        (a) $P(X>10)$   & (b) $P(X>20)$     & (c) $P(X<30)$ \\
        \multicolumn{3}{l}{(d) Find the value of $x$ such that $P(X<x)=0.90$.} \\
    \end{tabular}
}
\enanswer{
    For a random variable $X$ that is exponentially distributed with mean $10$ (i.e., parameter $\lambda = \frac{1}{10}$), we have that
    \[
        P(X > x) = e^{-\lambda x} = e^{-\frac{1}{10} x}
    \]
    We can calculate the desired probabilities given this formula:
    \begin{enumerate}[label=(\alph*)]
        \item $P(X > 10) = e^{-\frac{1}{10} \cdot 10} = e^{-1} \approx 0.3679$.
        \item $P(X > 20) = e^{-\frac{1}{10} \cdot 20} = e^{-2} \approx 0.1353$.
        \item $P(X < 30) = 1 - P(X > 30) = 1 - e^{-\frac{1}{10} \cdot 30} = 1 - e^{-3} \approx 0.9502$.
        \item $P(X < x) = 0.90$ means $P(X > x) = e^{-\frac{1}{10} \cdot x} = 0.10$.
        Therefore $-\frac{1}{10} \cdot x = \ln{0.10} = -\ln{10}\Rightarrow x = 10 \ln{10} \approx 23.0259$.  
    \end{enumerate}
}

\enquestion{4.103}{
    The lifetime of a mechanical assembly in a vibration test is exponentially distributed with a mean of $400$ hours.
    \begin{enumerate}[label=(\alph*)]
        \item What is the probability that an assembly on test fails in less than $200$ hours?
        \item What is the probability that an assembly operates for more than $500$ hours before failure?
        \item If an assembly has been on test for $400$ hours without a failure, what is the probability of a failure in the next $100$ hours?
        \item If $10$ assemblies are tested, what is the probability that at least one fails in less than $100$ hours? Assume that the assemblies fail independently.
        \item If $10$ assemblies are tested, what is the probability that all have failed by $800$ hours? Assume that the assemblies fail independently.
    \end{enumerate}
}
\enanswer{
    For a random variable $X$ that is exponentially distributed with mean $400$ (i.e., parameter $\lambda = \frac{1}{400}$), we have that
    \[
        P(X > x) = e^{-\lambda x} = e^{-\frac{1}{400} x}
    \]
    We can calculate the desired probabilities given this formula:
    \begin{enumerate}[label=(\alph*)]
        \item $P(X < 200) = 1 - P(X > 200) = 1 - e^{-\frac{1}{400} \cdot 200} = 1 - e^{-1/2} \approx 0.3935$.
        \item $P(X > 500) = e^{-\frac{1}{400} \cdot 500} = e^{-5/4} \approx 0.2865$.
        \item As the exponential distribution has the memoryless property, this probability is just
        \[
            P(X < 100) = 1 - P(X > 100) = 1 - e^{-\frac{1}{400} \cdot 100} = 1 - e^{-1/4} \approx 0.2212.
        \]
        \item Introduce a new variable counting the number of failing assemblies out of $10$ within $100$ hours.
        Then $Y$ is binomially distributed with $n = 10$ and $p = P(X < 100) = 0.2212$.
        The probability that at least one fails is Then
        \[
            P(Y \geq 1) = 1 - P(Y = 0) = 1 - \binompdf(n=10, p=0.2212, k=0) \approx 0.9179.
        \]
        \item Introduce a new variable counting the number of failing assemblies out of $10$ within $100$ hours.
        Then $Y$ is binomially distributed with $n = 10$ and $p = P(X < 800) = 1 - e^{-\frac{1}{400} \cdot 800} = 1 - e^{-2} \approx 0.8647$.
        The probability that all have failed by $800$ hours:
        \[
            P(Y = 10) = \binompdf(n=10, p=0.8647, k=10) \approx 0.2336.
        \]
    \end{enumerate}
}

\enquestion{4.165}{
    Suppose that $f(x) = 0.5x-1$ for $2<x<4$.
    Determine the following:

    \begin{tabular}{*{3}{m{.3\textwidth}}}
        (a) $P(X<2.5)$   & (b) $P(X>3)$     & (c) $P(2.5<X<3.5)$ \\
        \multicolumn{3}{l}{(d) Determine the cumulative distribution function of the random variable.} \\
        \multicolumn{3}{l}{(e) Determine the mean and variance of the random variable.} \\
    \end{tabular}    
}
\enanswer{
    We start with taking the antiderivative of $f(x)$ to obtain the cumulative distribution function (CDF):
    \begin{align*}
        P(X \leq x) = \int_{y=2}^x 0.5y-1 \dy &= \left[ 0.25y^2 - y \right|_{y=2}^x \\
                                                     &= \left(0.25x^2 - x\right) - \left( 1 - 2 \right) \\
                                                     &= 0.25x^2 - x + 1
    \end{align*}
    The cumulative distribution function (question (d)) $F(x)$ can then be given as
    \[
        F(x) = \begin{cases} 
            0, & \text{ for } x \leq 2 \\ 
            0.25x^2 - x + 1, & \text{ for } 2 < x < 4 \\ 
            1, & \text{ for } x \geq 4 \\ 
        \end{cases}
    \]
    We can use this CDF to compute the desired probabilities:
    \begin{enumerate}
        \item[(a)] $P(X < 2.5) = P(X \leq 2.5) = F(2.5) = 0.25 \cdot (2.5)^2 - 2.5 + 1 = 0.0625$
        \item[(b)] $P(X > 3) = 1 - P(X \leq 3) = 1 - F(3) = 1 - 0.25 \cdot 3^2 - 3 + 1 = 0.75$
        \item[(c)] 
        \begin{align*}
            P(2.5 < X < 3.5) &= P(X \leq 3.5) - P(X \leq 2.5) \\
            &= F(3.5) - F(2.5) \\ 
            &= \left(0.25 \cdot (3.5)^2 - 3.5 + 1\right) - \left(0.25 \cdot (2.5)^2 - 2.5 + 1\right)\\
            &= 0.5
        \end{align*}
        \item[(e)] We can compute the mean of this random variable as follows:
        \begin{align*}
        E(X) = \int_{x=-\infty}^{\infty} x \cdot f(x) \dx &= \int_{x=2}^{4} x \cdot (0.5x-1) \dx \\
                                                        &= \int_{x=2}^{4} \cdot 0.5x^2-x \dx \\
                                                        &= \left[ \frac{1}{6}x^3 - \frac{1}{2}x^2\right|_{x=2}^4 \\
                                                        &= \left(\frac{32}{3} - 8\right) - \left(\frac{4}{3} - 2\right) \\
                                                        &= \frac{10}{3}.
        \end{align*}
        For the variance, we first compute $E(X^2)$:
        \begin{align*}
            E(X^2) = \int_{x=-\infty}^{\infty} x^2 \cdot f(x) \dx &= \int_{x=2}^{4} x^2 \cdot (0.5x-1) \dx \\
                                                            &= \int_{x=2}^{4} 0.5x^3-x^2 \dx \\
                                                            &= \left[ \frac{1}{8}x^4 - \frac{1}{3}x^3\right|_{x=2}^4 \\
                                                            &= \left(32 - \frac{64}{3}\right) - \left(2 - \frac{8}{3}\right) \\
                                                            &= \frac{34}{3}.
        \end{align*}
        We now use the identity $V(X) = E(X^2) - (E(X))^2$:
        \begin{align*}
            V(X) = E(X^2) - (E(X))^2 = \frac{34}{3} - \left(\frac{10}{3}\right)^2 = \frac{2}{9}.
        \end{align*}
    \end{enumerate}
}